\chapter{估值与交易结构入门}

\section{可比公司与可比交易}

上市公司估值倍数（EV/EBITDA、P/E 等）提供市场锚点，但需调整：流动性折价、控制权溢价、行业周期位置与会计差异。\textbf{可比交易}更接近一级市场，但样本稀疏、披露不全，应对区间取保守情景。两者结合形成\textbf{估值走廊}，而非单点数字。

\section{DCF 与情景分析}

现金流折现（DCF）强迫显式列出增长、利润率、资本支出与终值假设；敏感性分析可展示关键驱动因子。早期企业 DCF 权重常低于\textbf{最近融资价格}或\textbf{风险投资法}，但仍有内部讨论价值。注意 WACC 与终值增长率假设勿机械套用教科书 defaults。

\section{股权比例与对价形式}

估值决定入门股权比例；\textbf{员工期权池}是否投前扩池会稀释有效比例。对价可为现金、换股或分期支付；分期与\textbf{earn-out} 绑定未来业绩，降低买方 upfront 风险但增加争议与会计复杂性。

\section{对赌与调整机制}

业绩承诺、回购与现金补偿等条款在中文语境常称「对赌」，实质是\textbf{价格与风险的再分配}。设计时应避免不可执行的惩罚、与公司法冲突的约定，以及税务上的意外后果。所有对赌须由\textbf{律师与税务顾问}联合审阅；本书仅强调经济含义：对赌不是「保证收益」，而是条件支付。

\section{谈判中的艺术}

估值是\textbf{区间博弈}：卖方要叙事上限，买方要风险下限。交易结构（分期、托管、陈述保证保险）常比硬压估值更能成交。保留\textbf{双赢叙事}有利于投后合作——过度压榨条款往往在董事会层面反噬。
